% !TeX program = xelatex
\documentclass[UTF8,12pt,a4paper]{ctexart}
\usepackage[margin=2cm,headheight=15pt]{geometry}
\usepackage{booktabs,tabularx,array,xcolor,hyperref,fancyhdr,tikz}
\usetikzlibrary{arrows.meta,positioning,fit}
\hypersetup{colorlinks=true,urlcolor=blue,linkcolor=black}
\pagestyle{fancy}\fancyhf{}\lhead{Python Programming}\rhead{Project 1}\cfoot{\thepage}
\setlength{\parskip}{0.35em}
\definecolor{workred}{HTML}{C73535}
\begin{document}
\begin{titlepage}\centering\vspace*{2cm}
{\Huge\bfseries 深圳大学\par}\vspace{1cm}
{\Large Python Programming\par}\vspace{0.8cm}
{\Huge\color{blue!65!black}课程实验报告\par}\vspace{0.5cm}
{\Large Project 1：基于 NumPy 的图像编码器前向计算\par}\vspace{2cm}
\begin{tabular}{ll}
课程名称：& Python Programming\\[0.5em]
实验项目：& NumPy Image Encoder\\[0.5em]
学生姓名：& \underline{\hspace{7cm}}\\[0.5em]
学号 / 专业班级：& \underline{\hspace{6cm}}\\[0.5em]
指导教师：& \underline{\hspace{7cm}}\\[0.5em]
实验日期：& \underline{\hspace{7cm}}\\[0.5em]
报告日期：& 2026 年 11 月 16 日
\end{tabular}
\vfill{\color{workred}\fbox{红色标记在正文中用于区分本次补全的代码部分。}}\end{titlepage}

\section*{摘要}
本实验使用 NumPy 完成 OpenCLIP ViT-B-32 图像编码器的前向计算。输入为给定预处理后的图像，输出为每张图像的 512 维单位向量。实现包括图像块切分、patch/CLS/位置嵌入、多头注意力、Transformer 残差与 MLP，以及 CLS 输出头。项目自带的完整数值检查全部通过；输出 token 最大绝对误差为 $5.439\times10^{-7}$。

\section{目的与环境}
目标是仅用 NumPy 重现已提供模型权重的图像前向过程，并按阶段对照独立 FP32 参考。运行环境：Windows、Python 3.14.0、NumPy 2.5.1、Pillow 12.3.0；CPU 运行。项目含 10 张图像、152 个视觉张量和参考数组。

\section{前向流程}
\begin{center}
\begin{tikzpicture}[font=\small, node distance=0.48cm, >={Stealth[length=2mm]},
provided/.style={draw=gray!65,fill=gray!8,rounded corners,minimum height=1.5cm,text width=1.7cm,align=center},
mine/.style={draw=workred,fill=workred!7,rounded corners,minimum height=1.5cm,text width=1.9cm,align=center,line width=1pt},
flow/.style={->,thick,gray!75}]
\node[provided] (in) {输入图像\\$[B,3,224,224]$\\预处理};
\node[mine,right=of in] (patch) {\textcolor{workred}{01 Patchify}\\$[B,49,3072]$};
\node[mine,right=of patch] (embed) {\textcolor{workred}{05 Embedding}\\patch + CLS + pos\\$[B,50,768]$};
\node[mine,right=of embed,text width=2.4cm] (block) {\textcolor{workred}{02--04 Transformer}\\12 blocks\\Attention + MLP};
\node[mine,right=of block] (out) {\textcolor{workred}{06 Output}\\CLS + projection\\$[B,512]$};
\draw[flow] (in)--(patch);\draw[flow] (patch)--(embed);\draw[flow] (embed)--(block);\draw[flow] (block)--(out);
\end{tikzpicture}
\end{center}
\begin{center}\color{workred}\small 红色模块为本次补全的六组代码。\end{center}

\section{本次补全的代码（红色高亮）}
\colorbox{workred!7}{\parbox{0.96\linewidth}{\color{workred}
TODO-01 Patchify；TODO-02/03 多头拆分与注意力；TODO-04 两条残差、Norm 2、GELU 与 MLP；TODO-05 patch/CLS/位置嵌入；TODO-06 CLS 输出头。\\
关键轴操作：$K\in[B,H,N,D_h]$，使用 \texttt{swapaxes(-1,-2)} 得到 $[B,H,D_h,N]$，保留 batch/head 轴。}}

\section{数值结果}
运行 \texttt{python check.py --stage all}。预处理 atol=$2\times10^{-6}$；中间张量 atol=$2\times10^{-4}$、rtol=$2\times10^{-4}$；输出向量 atol=$2\times10^{-5}$、rtol=$2\times10^{-4}$。
\begin{center}\small
\begin{tabular}{lrrr}\toprule
阶段 & 实际形状 & 最大绝对误差 & 平均绝对误差\\\midrule
preprocess &(10,3,224,224)&0.000e+00&0.000e+00\\
patches &(2,49,3072)&0.000e+00&0.000e+00\\
tokens &(2,50,768)&4.005e-05&1.289e-06\\
block0 &(2,50,768)&7.063e-06&3.250e-07\\
last\_tokens &(2,50,768)&4.768e-05&2.696e-06\\
output &(10,512)&5.439e-07&8.311e-08\\
batch &(3,512)&4.619e-07&7.632e-08\\\bottomrule
\end{tabular}\end{center}
全部阶段通过，patch 数据逐值一致；重复/重排 batch 检查通过。

\subsection{形状变化与问题修复}
patchify 把 $[B,3,224,224]$ 改排为 $[B,49,3072]$；注意力将 $[B,N,768]$ 拆分为 $[B,12,N,64]$。二维课堂示例中的 \texttt{k.T} 不能直接用于四维 key，因为它会反转所有轴。实现改用 \texttt{swapaxes(-1,-2)}，只交换末两轴，正确形成 $[B,H,64,N]$。

\section{复现步骤}
在 project1-numpy-student 目录运行：
\begin{verbatim}
python -m pip install -r requirements.txt
python check.py --stage all
python run.py --output output_tokens.npz
python run.py --images path/to/a.jpg path/to/b.png --output my_tokens.npz
\end{verbatim}
依赖为 NumPy、Pillow、imageio。默认输出保存图像 ID 和对应 $[B,512]$ float32 向量。

\section{贡献、协助与结论}
学生提供课程项目文件与提交要求，并确定实现/报告任务。Codex AI 协助补全六组 NumPy 实现、添加代码注释并起草本报告；数值验证由项目自带 check.py 实际执行。AI 协助已如实披露，提交前应由学生复核并确认理解。

本项目完成了从图像切块到归一化 CLS 向量的 NumPy 前向过程。独立 FP32 对照和 batch 顺序检查均通过。复现中最需要留意的是张量轴语义、线性权重方向与残差支路。

\section*{代码与模板链接}
项目代码 GitHub：\url{https://github.com/furinalover95/python-homework}\\
SZU 报告模板：\url{https://github.com/ChaoFan996/Experiment-Report-Template-SZU}

\vspace{0.5em}\noindent\textbf{提交文件：} image\_encoder.py、README.md、experiment\_report.pdf 和 experiment\_report.tex。报告正文 3 页，不含单独封面。
\end{document}
